\section{Trabajo futuro}\label{sec:conclusiones:futuro}

Sobre ACM2YAML quedan tres líneas de trabajo.
La primera es distinguir el límite publicado del aplicado, el elemento F-04, que reúne los límites que el proveedor declara y que la respuesta a la petición no refleja (Sección~\ref{sec:esquema:fuera}).
Es el caso de Mailchimp Transactional y de Amazon SES, que aceptaron llamadas por encima del tope que declaran.
Queda por decidir qué cifra va a la Datasheet y cómo se declara la fuente de cada una.
La segunda es fijar las doce formas esbozadas de la Tabla~\ref{tab:sinforma}, de una fecha de verificación para cada afirmación del registro de procedencia (F-09) al canal que ofrece cada operación (F-22), y contrastarlas con la guía de ACM2YAML.

La tercera es llevar como \textit{issues} las nueve propuestas de extensión de la Tabla~\ref{tab:propuestas}, cada una con su caso y el número de APIs en que aparece, al repositorio que tiene que cambiar: el de la guía si cambia el esquema y el de PRIME4API si el arreglo está en el motor (Sección~\ref{sec:esquema:metodo}).
De esas nueve, las cinco que van de los volúmenes contratados al plan con precio no público (P-01 a P-05) no figuran en ningún \textit{issue} de los repositorios de la guía, de PRIME4API ni de SPHERE (Sección~\ref{sec:esquema:fuera}).
Las cuatro restantes, del estado de la suscripción al comportamiento de la cuota al superarse (P-06 a P-09), aún deben contrastarse con los \textit{issues} abiertos antes de enviarse.
Las Datasheets del catálogo irían además al repositorio de la guía, en una \textit{pull request}, junto a las que ya publica \citep{apisconsumptionmodel}.

En las herramientas de consulta queda ejecutar con HARVEY4API las cuatro preguntas de clase H, las que combinan varias Datasheets u operaciones de PRIME4API, y comparar sus respuestas con las de Claude Opus 5.5 y con el cálculo a mano (Sección~\ref{sec:comparacion:harvey}).
También queda vincular la Datasheet del canal API con el iPricing, el metamodelo con el que SPHERE representa los planes y las \textit{features} de la página de precios \citep{garciafernandez2024pricing4saas}.
Con él, la pregunta de qué API da IP dedicada al menor precio para 200.000 correos al mes, que en la Sección~\ref{sec:comparacion:funcional} se calcula a mano, podría responderse con una consulta a SPHERE.

En SPHERE, la primera línea es implementar la sección de términos de consumo que la Sección~\ref{sec:propuesta:sphere:terminos} deja diseñada.
Lo publicado y lo medido entrarían por una ingesta que guarda la fecha de cada página, en un historial separado de las versiones de \textit{pricing} que SPHERE ya mantiene.
Ese diseño separa los límites propios de una interfaz, como las 10 conexiones SMTP concurrentes por IP que recomienda Postmark, de los compartidos entre varias.
Un límite compartido es la cuota diaria de MailerSend, que cuenta las peticiones de la API REST y las del relé SMTP.

Al diseño le falta una tercera clase, la de los límites globales, que valen para el servicio entero y no para una interfaz concreta.
Esa clase exige un nivel por encima de la interfaz.

La segunda línea en SPHERE es comparar las vías de contratación de un mismo servicio.
En RapidAPI, un mercado de APIs, hay fichas del servicio de envío de cuatro de las ocho APIs: SendGrid, Postmark, Brevo y Mailgun\footnote{Búsqueda sin cuenta en el buscador de RapidAPI, consultada el 23 de octubre de 2026.}.
Comparar el precio y los límites de SendGrid contratado en RapidAPI con los de su página de precios pide un \textit{pricing} por vía de contratación y el árbol de proveedor, servicio e interfaz de la maqueta de la colección para SPHERE (Sección~\ref{sec:propuesta:sphere:coleccion}).
Queda por aclarar si esas fichas, que publica un tercero, exigen una cuenta en el proveedor.

Las últimas líneas son \textit{skills}, instrucciones que sigue un agente basado en un modelo de lenguaje para leer los términos de consumo, y van de la más próxima a este trabajo a la más lejana.
Una agruparía las interfaces que comparten un límite, como la cuota diaria de MailerSend, y su salida se contrastaría con la página antes de entrar en la Datasheet.
Otra listaría como candidatas las contradicciones entre páginas de un mismo proveedor, como la guía de envío por lotes de Brevo, que admite 6.000 llamadas por hora al \textit{endpoint} de envío, 600 veces menos que los 3.600.000 por hora de su categoría General (Sección~\ref{sec:limites:publicado}).
Una tercera contaría los límites de cada versión fechada de una página, como las capturas de la página de límites de MailerSend.

La cuarta traduciría a cambios del dominio el diff sintáctico con el que SPHERE compara dos versiones de un pricing, que marca qué líneas difieren entre ellas \citep{sphere}.
La más lejana, porque depende de las otras, es publicarlas en un \textit{marketplace} que registre su uso real y conectarlas con SPHERE por el \textit{Model Context Protocol}, el mismo con el que HARVEY4API publica las operaciones de PRIME4API \citep{mcp}.

\phantomsection\label{sec:conclusiones:replicacion}Cualquiera de estas líneas parte del repositorio de datos del trabajo, \url{https://github.com/danvelcam/tfm-pricing-apis-correo}.
Reúne las ocho Datasheets con sus fuentes y evidencias junto a las herramientas del trabajo.
No incluye las credenciales de las APIs ni PRIME4API, y de las copias archivadas de las páginas de los proveedores publica solo la huella sha256.

Para comprobar la instalación, la herramienta \path{05_salida/_herramientas/verificar_prime4api.py} se ejecuta con \texttt{python3} sobre \path{03_datos/apis/postmark/datasheet.yaml}.
La primera vez clona PRIME4API en \path{~/.cache/atf/prime4api} sin fijar la revisión, así que antes hay que situar esa copia en la revisión \texttt{8e6011c} (Sección~\ref{sec:conclusiones:interna}).
El resultado debe coincidir con la fila de Postmark de la Tabla~\ref{tab:verificacion}: 11 comprobaciones superadas y 10 fallidas por valor ausente.
